\section{Discussion}\label{sec:discussion}

% DRAFT: to be completed after the field plots and the cost table are in, and after the
% full texts of the flagged related works have been read.

\paragraph{What the comparison measures}
By \cref{cor:training}, training on the energy is supervised regression in the stiffness
norm. The labels of the supervised networks carry no information that the assembled $K$
and $F$ do not, so the comparisons of \cref{sec:results2d,sec:results3d} are between two
losses: the stiffness-norm error, which can be evaluated without a solve, and the relative
displacement error, the usual supervised loss, which cannot. A supervised network trained
on the stiffness-norm error would share the label-free network's objective while paying for
the labels; we did not train one. The claim of this paper is therefore not that labels are
uninformative, but that the norm in which one can train without labels is the one that
controls the stresses, and that supervised training in the usual displacement norm does
not control them.

\paragraph{Why the label-free network's stresses are accurate}
By \cref{cor:training,prop:stress}, the label-free objective is, up to a constant, the
compliance-weighted mean-square error of the stress field. The supervised networks were
trained on the relative Euclidean error of the displacement, which weights all eigenmodes
of the stiffness matrix equally and so under-weights, relative to the energy norm, the
stiff modes that carry the strains (\cref{sec:theory:metrics}). A supervised network can
therefore reach a small displacement error with an error rich in stiff modes, and a large
stress error; instance No.~\numWorstIndex{} in seed 0 (\cref{sec:results3d:reverse}) is an
example, with a displacement error of \numWorstLabDisp{} and a relative energy gap of
\numWorstLabGap{}. The
energy term of the supervised transformer is a partial remedy: it added the stress-weighted
signal and removed most of the large failures, but on the same instances the label-free
transformer, which receives only that signal, was more accurate still.

\paragraph{The graph network}
The graph network had the lowest displacement error and among the largest stress errors.
On an instance where its displacement error is the lower and its energy-norm error the
higher, its error has the larger Rayleigh quotient $e^{\top}Ke/e^{\top}e$: relatively more
weight in stiff modes, as for an error that varies rapidly from node to node
(\cref{rem:rank}). We have not measured the spectral content of the errors; the field
plots of \cref{fig:field-worst,fig:field-median}, pending, will illustrate it. The observation matters for evaluation more than for the graph
network itself: a ranking of surrogates by displacement error alone can invert their
ranking by stress, the quantity on which design decisions usually depend.

\paragraph{The energy in deployment}
\Cref{cor:zero,prop:amp} give every surrogate, however it was trained, a check and a
repair that need only the assembled $K$ and $F$: the sign of $\Pi_h(u)$ flags a prediction
worse than the zero field, the difference of energies ranks the predictions of several
surrogates for the same instance exactly by their energy-norm errors, and the rescaling by
$\cstar$, which costs nothing once $\Pi_h(u)$ is known, removes any error of amplitude and
never increases the error in the energy norm. Every supervised network of
\cref{tab:2d,tab:3d} returned predictions that the check would have flagged; so, in
\numTwoFreeWorse{} of its \numTwoPairs{} instance-seed predictions, did the
two-dimensional label-free network.

\paragraph{Warm starts}
The label-free prediction did not shorten conjugate gradients to a relative residual of
\numKfiveTol{} in two dimensions (\cref{sec:results2d:decisions}), whereas the pretrain
finite element method has been reported to accelerate the warm-start stage of the solver by
up to an order of magnitude \cite{Wang2026PretrainFEM}, and neural operator warm starts
have been reported to accelerate iterative solvers \cite{Eshaghi2026NOWS}. \Cref{cor:warm}
suggests why the two kinds of result need not conflict: a warm start lowers the iteration
bound by a number of iterations set by its own accuracy, so the share of the run it can
save grows as the tolerance loosens and shrinks as it tightens, and a residual-based stop
can hide part of it (\cref{rem:residual}). Reported warm-start speed-ups are therefore
comparable only when the tolerance, the stopping test and the preconditioner are stated
with them. The comparison of solver settings with \cite{Wang2026PretrainFEM,Eshaghi2026NOWS}
is to be completed after their full texts have been read.

\paragraph{Transfer}
No network transferred to the finer mesh, and the readings of \cref{sec:transfer} point at
the output scaling: the network's output is multiplied by the largest nodal load, which
depends on the mesh, while the solution does not. The two-dimensional networks, whose
output was not scaled in this way, lost little accuracy on meshes \numCoarsenLow{} and
\numCoarsenHigh{} times finer, although that evidence is weak (\cref{sec:transfer}). A
separate pre-registered study tests an output scale independent of the mesh. Until it
reports, the label-free transformer of this paper should be regarded as valid within the
range of mesh sizes it was trained on.

\paragraph{Limitations}
The study is restricted to linear elastostatics, where \cref{lem:exact} is exact; one
family of geometries and one set of four load cases per dimension; and three seeds per
network. The two-dimensional runs of \cref{sec:results2d} used code that could not recover
the absolute load amplitude, and their supervised networks were not repeated with the later
code; their pre-registration documents were committed to the repository only after the runs
(\cref{sec:methods:prereg}). The scope of this paper, the energy objective without the
joint-embedding term that a pre-registered comparison kept, was chosen after that
comparison (\ref{app:prereg}). The
three-dimensional gate was amended after the initial run (\ref{app:prereg}), and both
readings are reported. The supervised baselines were trained on the relative displacement
error, the usual choice, and no supervised network was trained on the stiffness-norm error.
The reported von Mises error weights elements equally, so \cref{prop:stress} explains it
rather than bounds it. The inference cost (\cref{sec:cost:inference}), the field plots and
the per-load-case failure counts of the supervised networks are pending. No network
transferred to a substantially finer mesh.

\paragraph{Outlook}
The construction extends wherever each step of a computation is a convex minimisation. In
linear dynamics, every implicit step of a Newmark-type integrator solves a symmetric
positive definite system whose effective stiffness combines the stiffness, mass and
damping matrices, and \cref{lem:exact} holds step by step \cite{CaoSong2026}. Incremental
variational formulations of dissipative materials \cite{OrtizStainier1999} and variational
integrators \cite{MarsdenWest2001} write each step of a nonlinear history as the
minimisation of an incremental potential. There the energy remains a label-free training
objective, although the energy gap is no longer a squared norm, and the counterparts of
\cref{prop:amp,cor:zero} remain to be worked out. These are the problems in which a solve
may be expensive enough for label-free training to pay for itself.
